
\documentclass{article}
\usepackage{colortbl}
\usepackage{makecell}
\usepackage{multirow}
\usepackage{supertabular}

\begin{document}

\newcounter{utterance}

\centering \large Interaction Transcript for game `eqbench', experiment `eq\_bench', episode 104 with reflexion{-}Qwen3.5{-}4B.
\vspace{24pt}

{ \footnotesize  \setcounter{utterance}{1}
\setlength{\tabcolsep}{0pt}
\begin{supertabular}{c@{$\;$}|p{.15\linewidth}@{}p{.15\linewidth}p{.15\linewidth}p{.15\linewidth}p{.15\linewidth}p{.15\linewidth}}
    \# & \multicolumn{2}{c}{Player} && \multicolumn{2}{c}{Game Master} \\
    \hline

    \theutterance \stepcounter{utterance}  
    & & & \multicolumn{4}{p{0.6\linewidth}}{
        \cellcolor[rgb]{0.9,0.9,0.9}{
            \makecell[{{p{\linewidth}}}]{
                \texttt{\tiny{[P1$\langle$GM]}}
                \texttt{You have attempted this same task before and did not succeed. These are the plans you wrote after those attempts.} \\
\texttt{Trial \#1: New plan: I misinterpreted the emotional context of the dialogue. Terry's statement "Because plans save lives" is a defensive justification for his previous admission of having a plan, not a genuine belief that validates him against Cass's cynical attack. Cass's final line ("Or maybe they just make you feel better about yourself") is a direct challenge to Terry's ego and competence. Terry's response "You're goddamn right" is an aggressive, defensive retort meant to shut down the argument and reassert dominance, rather than a moment of vulnerability or vindication. Therefore, "Defensive" is the primary emotion, but "Vindictive" is too strong as it implies a desire for revenge, whereas Terry is simply protecting his ego. "Embarrassed" is incorrect because he is not admitting fault; he is doubling down. "Indifferent" is incorrect because the dialogue is highly charged. I need to ensure the scores reflect that the emotion is a reaction to being challenged on one's competence, making "Defensive" the highest score, while lowering "Vindictive" to reflect that it's a defensive posture rather than an active desire for harm. I will also ensure I strictly follow the output format without any extra text.} \\
\\ 
\texttt{Embarrassed: 0} \\
\texttt{Defensive: 9} \\
\texttt{Indifferent: 0} \\
\texttt{Vindictive: 1} \\
\\ 
\\ 
\\ 
\texttt{Your task is to predict the likely emotional responses of a character in this dialogue:} \\
\\ 
\texttt{Terry: We're not going in blind, Cass. Not this time.} \\
\texttt{Cass: And who's gonna stop me? You?} \\
\texttt{Terry: If I have to. } \\
\texttt{Cass: Oh, right. The big, bad Terry Holloway. Always with a plan. } \\
\texttt{Terry: You're goddamn right. You know why? Because plans save lives.} \\
\texttt{Cass: Or maybe they just make you feel better about yourself. Like you're in control.} \\
\texttt{{[}End dialogue{]}} \\
\\ 
\texttt{At the end of this dialogue, Terry would feel...} \\
\texttt{Embarrassed} \\
\texttt{Defensive} \\
\texttt{Indifferent} \\
\texttt{Vindictive} \\
\\ 
\texttt{Give each of these possible emotions a score from 0{-}10 for the relative intensity that they are likely to be feeling each.} \\
\\ 
\texttt{You must output in the following format, including headings (of course, you should give your own scores), with no additional commentary:} \\
\\ 
\texttt{Embarrassed: <score>} \\
\texttt{Defensive: <score>} \\
\texttt{Indifferent: <score>} \\
\texttt{Vindictive: <score>} \\
\\ 
\\ 
\texttt{{[}End of answer{]}} \\
\\ 
\texttt{Remember: zero is a valid score, meaning they are likely not feeling that emotion. You must score at least one emotion > 0.} \\
\\ 
\texttt{Your answer:} \\
            }
        }
    }
     \\ \\

    \theutterance \stepcounter{utterance}  
    & \multicolumn{4}{p{0.6\linewidth}}{
        \cellcolor[rgb]{0.561,0.737,0.561}{
            \makecell[{{p{\linewidth}}}]{
                \texttt{\tiny{[P1$\rangle$GM]}}
                \texttt{Embarrassed: 0} \\
\texttt{Defensive: 9} \\
\texttt{Indifferent: 0} \\
\texttt{Vindictive: 1} \\
            }
        }
    }
    & & \\ \\

    \theutterance \stepcounter{utterance}  
    & & & \multicolumn{2}{p{0.3\linewidth}}{
        \cellcolor[rgb]{0.9,0.9,0.9}{
            \makecell[{{p{\linewidth}}}]{
                \texttt{\tiny{[GM$|$GM]}}
                \texttt{Embarrassed: 0} \\
\texttt{Defensive: 9} \\
\texttt{Indifferent: 0} \\
\texttt{Vindictive: 1} \\
            }
        }
    }
    & & \\ \\

    \theutterance \stepcounter{utterance}  
    & & & \multicolumn{2}{p{0.3\linewidth}}{
        \cellcolor[rgb]{0.9,0.9,0.9}{
            \makecell[{{p{\linewidth}}}]{
                \texttt{\tiny{[GM$|$GM]}}
                \texttt{Embarrassed: 0} \\
\texttt{Defensive: 8} \\
\texttt{Indifferent: 0} \\
\texttt{Vindictive: 4} \\
            }
        }
    }
    & & \\ \\

    \theutterance \stepcounter{utterance}  
    & & & \multicolumn{2}{p{0.3\linewidth}}{
        \cellcolor[rgb]{0.9,0.9,0.9}{
            \makecell[{{p{\linewidth}}}]{
                \texttt{\tiny{[GM$|$GM]}}
                \texttt{game\_result = LOSE} \\
            }
        }
    }
    & & \\ \\

    \theutterance \stepcounter{utterance}  
    & & \multicolumn{4}{p{0.6\linewidth}}{
        \cellcolor[rgb]{0.95,0.95,0.95}{
            \makecell[{{p{\linewidth}}}]{
                \texttt{\tiny{[P1$\leftrightarrow$P]}}
                \texttt{Post{-}episode Reflexion (not sent to GM):} \\
\texttt{New plan: I misinterpreted the nuance of Terry's final response. While "Defensive" is correct, the score of 9 implies a high level of active resistance, whereas Terry's "You're goddamn right" is more of a weary, self{-}justifying assertion of his own logic rather than a purely reactive defense. More importantly, I underestimated the element of "Vindictive." Terry's entire monologue is built on the idea that his plans are superior and necessary ("Because plans save lives"), and he is actively dismissing Cass's cynical view as a delusion ("make you feel better"). This isn't just defense; it's a subtle attempt to assert moral superiority and dismiss the other person's perspective, which aligns with a low{-}level vindictive or smug attitude. I need to increase the "Vindictive" score to reflect this dismissive superiority while slightly lowering "Defensive" to reflect that it's a mix of stubbornness and self{-}justification rather than pure fear of being wrong. I will also ensure the output contains only the four lines with no introductory or concluding text.} \\
\\ 
\texttt{Embarrassed: 0} \\
\texttt{Defensive: 7} \\
\texttt{Indifferent: 0} \\
\texttt{Vindictive: 6} \\
            }
        }
    }
    & \\ \\

\end{supertabular}
}

\end{document}
